\section{智慧是我们和万物的关系}

本章收束到访谈后半程的核心定义：智慧就是我们和万物的关系。这个定义很特别，因为它没有把智慧理解为 IQ、知识量或推理分数，而是理解为人与人、人与物、人与组织、人与技术、人与世界的关系质量。AI 的意义也因此不只是替代人，而是把人从低价值、耗能的动作里解放出来，让人有更多能量去处理更高质量的关系。

\begin{figure}[H]
\centering
\includegraphics[width=0.92\textwidth]{wisdom-relation-map.png}
\caption{智慧是关系：智慧来自与人、物、组织、技术和世界持续接触。自制概念图，依据 02:26:47--02:31:10 对谈内容整理。}
\end{figure}

\begin{knowledgebox}{读图：智慧不是脑内分数}
人、物、组织、技术和世界共同构成智慧的来源。一个人如果只在脑内推理，而不接触真实人和真实事，智慧会变薄；一个 AI 如果只在 token 里循环，而不进入行动和反馈，也难以形成真实价值。
\end{knowledgebox}

\subsection{AI 的意义：减少消耗性动作}

上一段把智慧定义为关系，本节转向 AI 如何改变这些关系中的动作分配。李想举销售人员的例子：好的销售和产品专家喜欢与客户接触，因为那能带来能量；但重复邀约、低价值流程、机械沟通可能消耗能量。AI 的价值之一，是自动化这些消耗性动作，让人把能量放在更有智慧、更有关系质量的事情上。

\begin{importantbox}{AI 与人的关系}
AI 不只是替代劳动，更是重新分配人的能量。真正好的 AI 系统应该减少低价值、重复、耗能的动作，同时保留或增强人与客户、产品、组织和世界的高质量接触。
\end{importantbox}

\subsection{从套壳到基础能力}

如果 AI 的意义在于承担更多 action，那么最后就必须回到基础能力问题。在结尾处，李想谈到对“小红”等 Agent 应用的看法。他认为把 AI 往真正 action、专业 action 和 to-do list 方向推进，是重要尝试；但如果对手都是强模型公司，仅靠套壳、虚拟机和工具能力不够，必要的模型能力还是要真正训练。这段话给 AI 应用创业一个很清楚的提醒：入口、流程和产品体验可以先跑，但最终要补基础能力。

\begin{figure}[H]
\centering
\includegraphics[width=0.94\textwidth]{shell-vs-model-ability.png}
\caption{套壳与模型能力：应用创新要走向 action，最终仍要补基础能力。自制概念图，依据 02:44:21--02:45:30 对谈内容整理。}
\end{figure}

\begin{knowledgebox}{读图：套壳不是没价值，但不能停在那里}
左边的套壳/工具能快速验证用户需求和工作流，右边的基础能力决定长期上限。Agent 应用如果想进入专业 action，就必须逐步补模型、记忆、工具调用、评测和可靠性。
\end{knowledgebox}

\begin{warningbox}{不要把“套壳不是创新”理解成反应用}
应用创新非常重要，因为它最接近用户和行动场景。但当场景进入高价值、强对齐、高可靠任务时，只靠外部模型和薄工作流会遇到上限，基础能力必须跟上。
\end{warningbox}

\subsection{本章小结}

“智慧是关系”把整期从技术重新拉回人。VLA、Agent OS、世界模型和组织方法论最终都不是为了炫技，而是为了让智能更好地进入人与世界的关系中。

